\section{Physical Robot Experiments and Baseline Comparison}
\label{sec:results-baseline-comparison}

To evaluate GaitNet's continuous optimization strategy in
unstructured real-world scenarios, a baseline method is established
utilizing the \textit{Legged Perceptive Stack} \cite{leggedcontrol}.
This state-of-the-art framework employs Nonlinear Model Predictive
Control (NMPC) and Whole-Body Control (WBC) based on OCS2 \cite{leggedcontrol}.

In order to benchmark GaitNet and its specific structural ablations
against this baseline, a custom physical test environment was
constructed mirroring the simulation parameters
(\autoref{fig:figure-test-environment}). Hardware experiments are
conducted utilizing a Unitree Go1 quadruped. The physical terrain
consists of narrow strips of ground interspersed with gaps,
characterized by varying terrain difficulty metrics defined as the
percentage of unsteppable terrain (ranging from $0\%$ to $40\%$ difficulty).

We systematically evaluate the success rate, total steps taken, and
overall traversal time of the following methods across
\expBenchmarkEpisodes{} physical trials of \expBenchmarkDuration{} each:
\begin{enumerate}
  \item \textbf{Legged Perceptive Stack:} The model predictive
    control baseline \cite{leggedcontrol}.
  \item \textbf{GaitNet-Single:} A targeted ablation of GaitNet
    restricted to only one leg in swing at any given time.
  \item \textbf{GaitNet-Trot:} A targeted ablation of GaitNet
    restricted to only opposing legs in swing concurrently.
  \item \textbf{GaitNet (Full):} The proposed unlimited acyclic
    footprint generation model.
\end{enumerate}

An attempt is considered successful if the robot completes the
traversal without falling or stepping into the structurally removed
terrain sections (invalid terrain).

\begin{figure}[H]
  \centering
  \includegraphics[width=\textwidth]{images/figures/test-environment.png}
  % TODO: Replace with Real World test course image
  \caption{Unitree Go1 navigating the physical test environment used
    for baseline comparison. The terrain consists of narrow strips with
    missing sections mapping to varying terrain difficulties
  (percentage of unsteppable terrain).}
  \label{fig:figure-test-environment}
\end{figure}

\subsection{Physical Locomotion Results}

\autoref{tab:hardware-metrics} and
\autoref{fig:data-experiments-survival-curr} summarize the
performance of GaitNet (Full and ablations) alongside the baseline
method across a range of terrain difficulties spanning from flat
ground to $40\%$ unsteppable gaps.

\begin{table}[H]
  \centering
  \caption{Placeholder: Real-World Metrics on Unitree Go1}
  \begin{tabular}{lccc}
    \toprule
    \textbf{Method} & \textbf{Success Rate} & \textbf{Avg. Traversal
    Time [s]} & \textbf{Avg. Steps Taken} \\
    \midrule
    Legged Perceptive Stack & XX\% & XX.X & XXX \\
    GaitNet-Single & XX\% & XX.X & XXX \\
    GaitNet-Trot & XX\% & XX.X & XXX \\
    \textbf{GaitNet (Full)} & \textbf{XX\%} & \textbf{XX.X} & \textbf{XXX} \\
    \bottomrule
  \end{tabular}
  \label{tab:hardware-metrics}
\end{table}

\begin{figure}[H]
  \centering
  \includegraphics[width=\smallplotsize{}]{images/data/experiments/Gaitnet
  - survival-curr - 69.4 - individual.png} % TODO: Replace with new
  % hardware chart
  \caption{Hardware Evaluation Success Rates over terrain
    difficulties. GaitNet (Full) achieves superior fault-free traversal
    spanning challenging obstacle densities compared to restrictive
  gaits and the perceptive baseline.}
  \label{fig:data-experiments-survival-curr}
\end{figure}

\subsection{Discussion}

GaitNet's unconstrained performance relative to the Legged Perceptive
Stack and fixed-gait ablations reveals the clear advantage of
dynamic, acyclic spatial optimization. The system's ability to
coordinate multi-leg motions and adapt continuous spatial timing
concurrently minimizes the total steps required while avoiding
invalid terrain. This improvement is particularly pronounced at
higher terrain difficulties (e.g., $30$-$40\%$ gap densities), where
restrictive trot or single-leg strategies inherently clash with the
unstructured valid footprint areas.
